\section{Robots.txt Restrictions, Filing-Window Illiquidity, and Drift}
\label{sec:robots_txt_illiquidity_drift}

\subsection{Treatment Definition}
\label{subsec:robots_treatment}

I construct the treatment directly from archived \texttt{robots.txt}
rules.  For each firm-month, I identify whether at least one observed
corporate host restricts automated search-bot access through a
\texttt{Disallow} rule.  The main treatment is a search-bot-specific
restriction, denoted \textit{SearchBotSpec}.  It equals one once the firm
adopts a bot-specific \texttt{User-agent} rule that blocks at least one
search-oriented crawler, and zero before adoption.  I require treated
firms to have an observed clean pre-treatment month in which the relevant
search-bot group is not blocked, so firms already blocked at the beginning
of the robots.txt panel are excluded as left-censored.  Firms with valid
robots.txt observations that never block the relevant bots form the
clean-never control group.

For firm \(i\), let \(G_i^{SB}\) denote the first calendar month in which
a search-bot-specific block is observed.  For filing event \(e\) in month
\(m(e)\), the post-treatment indicator is
\[
    \text{Post}^{SB}_{i,e}
    =
    \mathbf{1}\{m(e)>G_i^{SB}\}.
\]
Thus the first blocking month itself is not treated as post.  The
filing-window regressions exclude filings in the first blocking month.
I also report robustness using broader robots.txt treatments: any
search-bot restriction, any bot-specific restriction across the broader
crawler set, and full-site bot-specific restrictions.

\begin{table}[!htbp]
\centering
\caption{Robots.txt Treatment Definitions}
\label{tab:robots_treatment_definitions}
\begin{threeparttable}
\begin{tabular}{llp{7.5cm}}
\toprule
Variable & Label & Definition \\
\midrule
\textit{SearchBotSpec} & Main treatment &
Firm adopts a bot-specific \texttt{User-agent} rule that blocks at least
one search-oriented crawler. \\
\textit{BotSpec} & Bot-specific robustness &
Firm adopts a bot-specific \texttt{User-agent} rule that blocks at least
one crawler in the broader tracked crawler set. \\
\textit{RobotsAny} & Broad robots.txt restriction &
Firm blocks at least one tracked crawler through any robots.txt rule,
including wildcard rules. \\
\textit{RootSpec} & Full-site bot-specific restriction &
Firm adopts a bot-specific rule with \texttt{Disallow: /}, restricting
the full site for the targeted crawler. \\
\bottomrule
\end{tabular}
\begin{tablenotes}[flushleft]\footnotesize
\item Treated firms must exhibit an observed clean-to-block transition
within the robots.txt panel.  Left-censored and future-treated firms are
excluded.  Clean-never firms are retained as controls.
\end{tablenotes}
\end{threeparttable}
\end{table}

\subsection{Filing-Window Illiquidity}
\label{subsec:robots_amihud}

I measure filing-window illiquidity using Amihud's price-impact measure,
\[
    \text{Amihud}_{i,d}
    =
    10^6 \times \frac{|R_{i,d}|}{\text{DollarVolume}_{i,d}},
\]
following \citet{Amihud2002}.  For each filing window, the outcome is the
average filing-window Amihud measure minus the firm's pre-filing benchmark
average over days \([-30,-6]\).  The main event windows are \([0,+1]\),
\([0,+3]\), \([0,+5]\), \([-5,+5]\), \([+6,+30]\), and \([0,+30]\).  The
estimating equation absorbs firm and filing-month fixed effects.

\begin{table}[!htbp]
\centering
\caption{Robots.txt Restrictions and Filing-Window Amihud Illiquidity}
\label{tab:robots_amihud_main}
\begin{threeparttable}
\begin{tabular}{lcccccccc}
\toprule
Treatment &
\multicolumn{3}{c}{BJS, \([0,+5]\)} &
\multicolumn{5}{c}{TWFE companion, \([0,+5]\)} \\
\cmidrule(lr){2-4}\cmidrule(lr){5-9}
& Est. & \(p\) & Pretrend \(p\)
& Est. & \(p\) & \(N\) & Firms & Within \(R^2\) \\
\midrule
SearchBotSpec & 0.054 & 0.006 & 0.535 & 0.057 & 0.049 & 69,529 & 4,351 & 0.0001 \\
BotSpec       & 0.045 & 0.001 & 0.416 & 0.046 & 0.033 & 67,693 & 4,242 & 0.0001 \\
RobotsAny     & 0.025 & 0.110 & 0.985 & 0.019 & 0.213 & 30,881 & 1,947 & 0.0001 \\
RootSpec      & 0.050 & 0.007 & 0.152 & 0.050 & 0.096 & 69,224 & 4,333 & 0.0000 \\
\bottomrule
\end{tabular}
\begin{tablenotes}[flushleft]\footnotesize
\item The dependent variable is abnormal average Amihud illiquidity in the
filing window.  BJS estimates use imputation DID with firm clustering and
12-month pretrend tests.  TWFE companion estimates include firm and
filing-month fixed effects with firm-clustered standard errors.
\end{tablenotes}
\end{threeparttable}
\end{table}

The search-bot-specific treatment produces a positive and statistically
significant filing-window increase in illiquidity.  The effect is
concentrated around the filing date: the \([0,+5]\) estimate is large and
significant, while delayed windows such as \([+6,+30]\) and \([0,+30]\)
are economically small and statistically insignificant.  The broader
\textit{RobotsAny} treatment is attenuated, consistent with wildcard and
general crawl restrictions pooling together more heterogeneous firm
policies.

\begin{table}[!htbp]
\centering
\caption{Pre-Filing Benchmark Robustness for Abnormal Amihud}
\label{tab:amihud_benchmark_robustness}
\begin{threeparttable}
\begin{tabular}{lccc}
\toprule
Pre-filing benchmark & \([0,+5]\) estimate & \(p\)-value & \([+6,+30]\) \(p\)-value \\
\midrule
\([-30,-6]\)  & 0.0369 & 0.005 & 0.720 \\
\([-45,-6]\)  & 0.0353 & 0.009 & 0.415 \\
\([-60,-6]\)  & 0.0374 & 0.017 & 0.218 \\
\([-60,-16]\) & 0.0303 & 0.047 & 0.596 \\
\bottomrule
\end{tabular}
\begin{tablenotes}[flushleft]\footnotesize
\item The table varies the pre-filing benchmark used to construct abnormal
Amihud.  The delayed-window placebo remains insignificant across all
benchmark definitions.
\end{tablenotes}
\end{threeparttable}
\end{table}

\subsection{Post-Earnings-Announcement Drift}
\label{subsec:robots_pead}

I also test whether robots.txt restrictions amplify
post-earnings-announcement drift.  This test follows the PEAD literature,
using announcement and drift windows similar to \citet{DellaVignaPollet2009}
and \citet{HirshleiferLimTeoh2009}, and explicitly accounting for the
post-2000 decline in PEAD documented by \citet{Martineau2021}.  The drift
outcome is benchmark-adjusted buy-and-hold abnormal return over
post-announcement windows.  Surprise ranks are constructed from either
random-walk seasonal earnings surprise or analyst forecast error.  I first
verify that baseline PEAD remains detectable in the sample; consistent
with \citet{Martineau2021}, the remaining signal is concentrated in
microcap firms and short post-announcement windows.

I then estimate whether the PEAD slope changes after robots.txt treatment:
\[
    BHAR_{i,q,w}
    =
    \beta\,\text{SurpriseRank}_{i,q}
    +
    \alpha\,\text{PostTreat}_{i,q}
    +
    \gamma\,
    \text{SurpriseRank}_{i,q}\times \text{PostTreat}_{i,q}
    +
    \mu_i+\tau_q+\varepsilon_{i,q,w}.
\]
The coefficient of interest is \(\gamma\).  A positive \(\gamma\) would
imply that robots.txt restrictions amplify drift by steepening the
relation between earnings surprise and post-announcement abnormal returns.

\begin{table}[!htbp]
\centering
\caption{Robots.txt Restrictions and PEAD Amplification}
\label{tab:robots_pead_drift}
\begin{threeparttable}
\begin{tabular}{llccccc}
\toprule
Treatment & Cell & Implied \(D10-D1\) effect & \(p\) & \(N\) & Treated firms & Within \(R^2\) \\
\midrule
SearchBotSpec & SUE microcap, \([+2,+5]\)    &  5.26\% & 0.360 & 31,417 &  93 & 0.0024 \\
SearchBotSpec & SUE microcap, \([+2,+60]\)   &  0.69\% & 0.887 & 31,418 &  93 & 0.0007 \\
SearchBotSpec & Analyst FE full, \([+2,+60]\) & -1.98\% & 0.712 & 36,041 &  79 & 0.0020 \\
BotSpec       & SUE microcap, \([+2,+5]\)    &  5.41\% & 0.287 & 30,733 & 116 & 0.0025 \\
BotSpec       & SUE microcap, \([+2,+60]\)   & -1.28\% & 0.785 & 30,734 & 116 & 0.0007 \\
RobotsAny     & SUE microcap, \([+2,+5]\)    &  1.44\% & 0.059 & 14,455 & 416 & 0.0037 \\
RobotsAny     & SUE microcap, \([+2,+60]\)   &  2.83\% & 0.248 & 14,456 & 416 & 0.0006 \\
Broad discovery & SUE microcap, \([+2,+60]\) & -1.39\% & 0.706 & 31,029 & 144 & 0.0007 \\
Broad discovery & Analyst FE full, \([+2,+60]\) & -2.82\% & 0.226 & 33,882 & 196 & 0.0022 \\
\bottomrule
\end{tabular}
\begin{tablenotes}[flushleft]\footnotesize
\item Each estimate reports \(\gamma \times 9\), the implied change in the
\(D10-D1\) drift spread after treatment.  All regressions include firm and
fiscal-quarter fixed effects and use firm-clustered standard errors.
Two-way firm-by-quarter clustering gives identical point estimates; singular
two-way cells are flagged separately in the appendix.
\end{tablenotes}
\end{threeparttable}
\end{table}

The drift results do not support a delayed-earnings-news channel.  Even in
the microcap cells where baseline PEAD remains alive, robots.txt
restrictions do not robustly amplify post-announcement drift.  The only
marginal estimate appears for the broadest \textit{RobotsAny} treatment in
the short \([+2,+5]\) window, and it does not persist over the cumulative
\([+2,+60]\) drift window.  For the broad discovery treatment that produces
the filing-window Amihud effect, the drift-window interactions are null or
negative.  These findings suggest that the main effect operates through a
filing-window liquidity or price-impact channel, rather than through
delayed earnings-news diffusion.

